%! Symmetric Positive Definite Matrices — Unit 6, Lesson 6.3 — Homework
\documentclass[10pt]{article}
\usepackage{linalg-article}
\usepackage{linalg-boxes}
\newcommand{\TallMath}[1]{$\displaystyle #1\rule[-1.4em]{0pt}{3.2em}$}
\newcommand{\xx}{\mathbf{x}}
\newcommand{\zero}{\mathbf{0}}
\newcommand{\T}{\mathsf{T}}

\begin{document}

\pageheader{Unit 6, Lesson 6.3}{Homework}
\namedateperiod

\vspace{0.12in}

\begin{notesbox}{Practice}
A \emph{symmetric} matrix ($S=S^{\T}$) has perpendicular eigenvectors, so normalizing them gives an
orthonormal $Q$ and $S=Q\Lambda Q^{\T}$. It is \emph{positive definite} when every eigenvalue is
positive --- equivalently $a>0$ and $ac-b^2>0$.

\begin{enumerate}[leftmargin=1.9em, itemsep=12pt]
  \item \textbf{Spectral factorization.} For $S=\begin{bmatrix} 2 & 1 \\ 1 & 2 \end{bmatrix}$ (eigenvalues
        $\lambda=3,1$ with eigenvectors $\begin{bsmallmatrix}1\\1\end{bsmallmatrix},\begin{bsmallmatrix}1\\-1\end{bsmallmatrix}$),
        normalize the eigenvectors into an orthonormal $Q$, write $\Lambda$, and state $S=Q\Lambda Q^{\T}$. \writelines{3}
  \item \textbf{Positive-definite screen.} Each matrix is symmetric. Use the quick test ($a>0$ and $ac-b^2>0$)
        to decide whether it is positive definite.
    \begin{enumerate}[label=(\alph*), leftmargin=1.8em, itemsep=4pt]
      \item $\begin{bmatrix} 2 & 1 \\ 1 & 2 \end{bmatrix}$ \hfill
      \item $\begin{bmatrix} 1 & 3 \\ 3 & 1 \end{bmatrix}$ \hfill
      \item $\begin{bmatrix} 1 & 2 \\ 2 & 4 \end{bmatrix}$
    \end{enumerate}
    \vspace{1.3cm}
  \item \textbf{Energy.} For $S=\begin{bmatrix} 2 & 1 \\ 1 & 2 \end{bmatrix}$, compute the energy $\xx^{\T}S\xx$
        at $\xx=\begin{bsmallmatrix}1\\0\end{bsmallmatrix}$ and at $\xx=\begin{bsmallmatrix}1\\1\end{bsmallmatrix}$.
        How do the results fit with $S$ being positive definite? \writelines{2}
  \item \textbf{Positive definite vs.\ invertible.} A symmetric matrix has eigenvalues $\lambda=6$ and
        $\lambda=-2$. Is it positive definite? Is it invertible? Explain how these two questions differ. \writelines{2}
  \item \textbf{Explain.} (a) Why does normalizing perpendicular eigenvectors make $Q^{-1}=Q^{\T}$?
        (b) Why must a positive-definite matrix be invertible? \writelines{2}
\end{enumerate}
\end{notesbox}

\begin{extensionbox}
\textbf{The geometry.} For a positive-definite $S$, the set of vectors with energy $\xx^{\T}S\xx=1$ traces an
\emph{ellipse}, and its axes point along the eigenvectors with lengths $1/\sqrt{\lambda}$. For
$S=\begin{bsmallmatrix}2&1\\1&2\end{bsmallmatrix}$ ($\lambda=3,1$), along which eigenvector direction is the
ellipse \emph{longest}, and why? \writelines{2}
\end{extensionbox}

\begin{spiralbox}
\textbf{Coming next --- Lesson 6.4: Systems of Differential Equations.} You have now met eigenvalues,
diagonalization ($A=X\Lambda X^{-1}$), and the symmetric case ($S=Q\Lambda Q^{\T}$). Next you'll put them to
work on a system that \emph{changes over time} --- the eigenvalues tell you how each piece grows or decays.
\end{spiralbox}

\end{document}
